% Bagian: second-order-logic
% Bab: sol-and-sets
% Subbagian: power-of-continuum

\documentclass[../../../include/open-logic-section]{subfiles}

\begin{document}

\olfileid[id]{sol}{set}{pow}

\olsection{Kuasa Kontinuum}

\begin{explain}
Dalam logika orde kedua kita dapat menguantifikasi atas himpunan-himpunan
bagian !!{domain}, tetapi tidak atas himpunan yang terdiri atas
himpunan-himpunan bagian !!{domain}. Untuk melakukannya secara langsung,
kita memerlukan logika \emph{orde ketiga}. Misalnya, jika kita ingin
menyatakan Teorema Cantor bahwa tidak terdapat fungsi !!{injective} dari
himpunan kuasa suatu himpunan ke himpunan itu sendiri, kita dapat mencoba
merumuskannya sebagai ``untuk setiap himpunan~$X$ dan setiap himpunan~$P$,
jika $P$ adalah himpunan kuasa dari~$X$, maka tidak berlaku
$\cardle{P}{X}$''. Untuk mengatakan bahwa $P$ adalah himpunan kuasa dari~$X$,
kita perlu memformalkan bahwa !!{element}s dari $P$ tepat merupakan semua
himpunan bagian dari~$X$, sehingga diperlukan sesuatu seperti
$\lforall[Y][(P(Y) \liff Y \subseteq X)]$. Masalahnya terletak pada $P(Y)$:
itu bukan !!a{formula} logika orde kedua, karena hanya suku yang dapat menjadi
argumen bagi variabel relasi satu-tempat seperti~$P$.

Namun, kita dapat \emph{menyimulasikan} kuantifikasi atas himpunan yang
beranggotakan himpunan jika domain cukup besar. Gagasannya ialah memanfaatkan
fakta bahwa
relasi dua-tempat~$R$ menghubungkan !!{element}s dari !!{domain} dengan
!!{element}s dari !!{domain}. Jika diberikan relasi~$R$ semacam itu, kita
dapat menghimpun semua !!{element} yang dengannya suatu~$x$ berelasi
menurut~$R$:
$\Setabs{y \in \Domain{M}}{R(x,y)}$ adalah himpunan yang ``dikodekan
oleh''~$x$. Sebaliknya, jika $Z \subseteq \Pow{\Domain{M}}$ adalah suatu
koleksi himpunan bagian dari~$\Domain{M}$ dan !!{element}s dalam
$\Domain{M}$ sekurang-kurangnya sama banyak dengan himpunan-himpunan
dalam~$Z$, maka terdapat pula relasi~$R \subseteq \Domain{M}^2$ sedemikian
sehingga setiap $Y \in Z$ dikodekan oleh suatu~$x$ melalui~$R$.
\end{explain}

\begin{defn}
Jika $R \subseteq \Domain{M}^2$, maka \emph{$x$ mengodekan dengan~$R$}
$\Setabs{y \in \Domain{M}}{R(x, y)}$.
\end{defn}

Jika !!a{element}~$x \in \Domain{M}$ mengodekan dengan~$R$ suatu himpunan
$Z \subseteq \Domain{M}$, maka suatu himpunan $Y \subseteq \Domain{M}$
mengodekan himpunan yang terdiri atas himpunan-himpunan, yakni
himpunan-himpunan yang dikodekan oleh !!{element}s dari~$Y$. Jadi, himpunan
$Y$ dapat mengodekan dengan~$R$ himpunan $\Pow{X}$. Hal itu terjadi jika dan
hanya jika untuk setiap $Z \subseteq X$, suatu $x \in Y$ mengodekan
dengan~$R$ himpunan~$Z$, dan setiap $x \in Y$ mengodekan dengan~$R$ suatu
himpunan~$Z \subseteq X$.

\begin{prop}
!!^{formula}
  \[
  \fn{Codes}(x, R, Z) \ident \lforall[y][(Z(y) \liff
    R(x, y))]
  \]
menyatakan bahwa $s(x)$ mengodekan dengan~$s(R)$ himpunan $s(Z)$.
!!^{formula}
\begin{multline*}
  \fn{Pow}(Y, R, X) \ident \\
  \lforall[Z][(Z \subseteq X \lif \lexists[x][(Y(x) \land
    \fn{Codes}(x, R, Z))])] \land {} \\ \lforall[x][(Y(x) \lif
  \lforall[Z][(\fn{Codes}(x, R, Z) \lif Z \subseteq X)])]
\end{multline*}
menyatakan bahwa $s(Y)$ mengodekan dengan~$s(R)$ himpunan kuasa
dari~$s(X)$, yakni anggota-anggota $s(Y)$ mengodekan dengan~$s(R)$ tepat
himpunan-himpunan bagian dari $s(X)$.
\end{prop}

\begin{explain}
Dengan siasat ini, kita dapat mengungkapkan pernyataan tentang himpunan kuasa
dengan menguantifikasi atas kode-kode himpunan bagian alih-alih atas himpunan
bagian itu sendiri. Misalnya, Teorema Cantor kini dapat diungkapkan dengan
mengatakan bahwa tidak terdapat fungsi !!{injective} dari domain sebarang
relasi yang mengodekan himpunan kuasa dari~$X$ ke~$X$ sendiri.
\end{explain}

\begin{prop}
Kalimat
\begin{multline*}
  \lforall[X][\lforall[Y][\lforall[R][(\fn{Pow}(Y, R, X) \lif \\
      \lnot \lexists[u][(\lforall[x][\lforall[y][((Y(x) \land Y(y) \land
            \eq[u(x)][u(y)]) \lif \eq[x][y])]] \land {}\\
        \lforall[x][(Y(x) \lif X(u(x)))])])]]]
\end{multline*}
valid.
\end{prop}

\begin{explain}
Himpunan kuasa dari himpunan !!a{denumerable} adalah !!{nonenumerable},
sehingga kardinalitasnya lebih besar daripada kardinalitas sebarang himpunan
!!{denumerable} (yaitu~$\aleph_0$). Ukuran $\Pow{\Nat}$ disebut ``kuasa
dari kontinuum,'' karena sama dengan ukuran himpunan titik-titik pada garis
bilangan real,~$\Real$. Jika !!{domain} cukup besar untuk mengodekan himpunan
kuasa dari suatu himpunan !!a{denumerable}, kita dapat mengungkapkan bahwa suatu
himpunan berukuran sama dengan kontinuum dengan mengatakan bahwa himpunan itu
berkardinalitas sama dengan sebarang himpunan~$Y$ yang mengodekan himpunan
kuasa dari suatu himpunan~$X$ berukuran~$\aleph_0$. (Jika domain tidak cukup
besar, yakni tidak memuat himpunan bagian yang berkardinalitas sama
dengan~$\Real$, maka tidak mungkin pula ada relasi yang
mengodekan~$\Pow{X}$.)
\end{explain}

\begin{prop}
Jika $\cardle{\Real}{\Domain{M}}$, maka !!{formula}
\begin{multline*}
\fn{Cont}(Y) \ident 
\lexists[X][\lexists[R][((\fn{Aleph}_0(X) \land
      \fn{Pow}(Y, R, X)) \land \\ 
      \lforall[x][\lforall[y][((Y(x) \land {} 
      Y(y) \land \lforall[z][R(x,z) \liff R(y,z)]) \lif \eq[x][y])]])]]
\end{multline*}
menyatakan bahwa $\cardeq{s(Y)}{\Real}$.
\end{prop}

\begin{proof}
  $\fn{Pow}(Y, R, X)$ menyatakan bahwa $s(Y)$ mengodekan dengan~$s(R)$
  himpunan kuasa dari~$s(X)$, yang menurut $\fn{Aleph}_0(X)$ terhitung.
  Jadi, $s(Y)$ sekurang-kurangnya sebesar kuasa kontinuum,
  meskipun mungkin lebih besar (jika beberapa anggota dari $s(Y)$
  mengodekan himpunan bagian yang sama dari~$X$). Kemungkinan ini disingkirkan
  oleh konjung terakhir, yang mensyaratkan bahwa pengaitan antara
  anggota-anggota $s(Y)$ dan himpunan-himpunan bagian dari $s(X)$ melalui
  $s(R)$ bersifat !!{injective}.
\end{proof}

\begin{prop}
$\cardeq{\Domain{M}}{\Real}$ jika dan hanya jika
\begin{multline*}
  \Sat{M}{\lexists[X][\lexists[Y][\lexists[R][(\fn{Aleph}_0(X) \land \fn{Pow}(Y, R, X) \land \\
          \lexists[u][(\lforall[x][\lforall[y][(\eq[u(x)][u(y)] \lif \eq[x][y])]] \land {} \\\lforall[y][(Y(y) \liff \lexists[x][\eq[y][u(x)]])])])]]]}.
        \end{multline*}
  \end{prop}

\begin{explain}
Hipotesis Kontinuum adalah pernyataan bahwa ukuran kontinuum merupakan
kardinalitas !!{nonenumerable} pertama, yakni bahwa $\Pow{\Nat}$
berukuran~$\aleph_1$.
\end{explain}

\begin{prop}
Hipotesis Kontinuum benar jika dan hanya jika \[\fn{CH} \ident
\lforall[X][(\fn{Aleph}_1(X) \liff \fn{Cont}(X))]\] valid.
\end{prop}

Perhatikan bahwa tidak benar bahwa $\lnot \fn{CH}$ valid jika dan hanya jika
Hipotesis Kontinuum salah. Dalam domain yang !!a{enumerable}, tidak terdapat
himpunan bagian berukuran~$\aleph_1$ dan juga tidak terdapat himpunan bagian
seukuran kontinuum, sehingga $\fn{CH}$ selalu benar dalam domain yang
!!a{enumerable}. Namun, kita dapat memberikan kalimat lain yang valid jika dan
hanya jika Hipotesis Kontinuum salah:

\begin{prop}
Hipotesis Kontinuum salah jika dan hanya jika \[\fn{NCH} \ident
\lforall[X][(\fn{Cont}(X) \lif \lexists[Y][(Y \subseteq X \land \lnot
    \fn{Count}(Y) \land \lnot \cardeq{X}{Y})])]\] valid.
\end{prop}

\end{document}
